% Tabla del cuerpo (Tabla de contrastes) + macros de estado (edición ES).
% La tabla de resultados formales vive en el apéndice del fichero principal.
\newcommand{\stproven}{\textcolor{sepblue}{\ding{51}}\,probado}
\newcommand{\stfuture}{\textcolor{restgrey}{$\circ$}\,futuro}
\newcommand{\stdef}{\textcolor{sepblue}{$\bullet$}\,definido}

% ---------- Tabla de contrastes computacionales ----------
\begin{table}[t]
  \centering
  \small
  \arrayrulecolor{sepblue}
  \caption{Contrastes computacionales (GAP / Sage, exactos). \emph{Izquierda:} la
  dicotomía de separación, verificada sin contraejemplo en todo grupo transitivo de
  grado pequeño, más el testigo de canon en $\Z/12$. \emph{Derecha:} una
  recomputación independiente de la tabla distinguidora de Chan
  $D(\GL_n(\mathbb{F}_q)\curvearrowright\mathbb{F}_q^n)$ para $n>2$, $q\le n{+}1$,
  coincidente con Klavžar--Wong--Zhu; el color de celda codifica $D$ (pálido $D{=}2$
  $\to$ intenso $D{=}4$), así que las dos excepciones destacan de un vistazo.}
  \label{tab:checks}
  \renewcommand{\arraystretch}{1.2}
  \begin{minipage}[t]{0.46\textwidth}
    \centering
    \begin{tabular}{@{}ll@{}}
      \toprule
      \rowcolor{sepblue}
      \textcolor{white}{\textbf{Contraste}} & \textcolor{white}{\textbf{Resultado}} \\
      \midrule
      \rowcolor{sepband}
      barrido dicotomía & \multirow{1}{*}{0 contraejemplos} \\
      \small 120 grps.\ transitivos, gr.\ 2--9 & \\
      canon $\Z/12$ $\{0,3,6,9\}{\oplus}\{0,1,2\}$ & tesela (extremal) \\
      \bottomrule
    \end{tabular}
  \end{minipage}\hfill
  \begin{minipage}[t]{0.5\textwidth}
    \centering
    \setlength{\tabcolsep}{5pt}
    \begin{tabular}{@{}c|cccc@{}}
      \multicolumn{1}{@{}c}{\scriptsize $n\backslash q$} & \scriptsize 2 & \scriptsize 3 & \scriptsize 4 & \scriptsize 5 \\
      \hline
      \scriptsize 3 & \cellcolor{dFour}\textbf{4} & \cellcolor{dTwo}2 & \cellcolor{dTwo}2 & --- \\
      \scriptsize 4 & \cellcolor{dThree}\textbf{3} & \cellcolor{dTwo}2 & \cellcolor{dTwo}2 & \cellcolor{dTwo}2 \\
      \scriptsize 5 & \cellcolor{dTwo}2 & \cellcolor{dTwo}2 & \cellcolor{dTwo}2 & \cellcolor{dTwo}2 \\
      \scriptsize 6 & \cellcolor{dTwo}2 & \cellcolor{dTwo}2 & \cellcolor{dTwo}2 & \cellcolor{dTwo}2 \\
    \end{tabular}
    \\[3pt]
    {\scriptsize\colorbox{dTwo}{\,$D{=}2$\,}\ \colorbox{dThree}{\,$D{=}3$\,}\ \colorbox{dFour}{\textcolor{white}{\,$D{=}4$\,}}\quad(``---'': $q$ no es potencia de primo)}
  \end{minipage}
\end{table}
